\chapter{Frames, Clocks, and the Interval}
\label{chap:frames-clocks-interval}

Part V turns to relativity and gravitation, where the record's own geometry bends. It opens with the structure
relativity rests on: no preferred frame, an isotropic metric, and a single invariant --- the interval, which
is the proper time a clock counts along its own path. Which frame we assign is ours to set --- none of them
privileged --- and the interval the world keeps the same in every one of them is the world's. Seven effects
carry the chapter, and one of them, the relativity of frames, splits into two: the relativity of simultaneity
and the maximality of proper time, two aspects distinct enough to take a section each. The count, which through four parts was carried, balanced, and
charged, here acquires a geometry --- and the geometry is the measure of the count itself.

\section{The Maxwell Effect: no preferred frame}
\label{se:maxwell-relativity}

A frame is a local labeling of refinement events, and no labeling is privileged --- the frame ours to assign,
none of them the world's to prefer. The fact that fixes this is the speed of light: Maxwell's equations give one
speed $c$ for light, the same in every frame, and a constant
speed forces a geometry in which a particular combination of time and space, not either alone, is invariant.
That combination is the interval. Between two nearby events separated by a time $dt$ and a distance $dx$, the
proper time elapsed is
\begin{equation}
  d\tau^2 = dt^2 - \frac{dx^2}{c^2},
\end{equation}
the Minkowski interval --- written in general with the metric as
\begin{equation}
  s^2 = g_{\mu\nu}\, dx^\mu dx^\nu,
\end{equation}
the metric $g_{\mu\nu}$ carrying the signs that set time against space. A moving clock trades time for distance:
the faster it goes, the more of its separation is spent as $dx$ and the less proper time $d\tau$ it
accumulates --- but the interval combining the two is the same for every observer: the split into time and
space ours to set with the frame, the interval the world holds fixed beneath it.

A transformation between frames is admissible exactly when it preserves this interval. The Lorentz
transformation,
\begin{equation}
  t' = \gamma\Bigl(t - \frac{vx}{c^2}\Bigr), \qquad x' = \gamma\,(x - vt), \qquad
  \gamma = \frac{1}{\sqrt{1 - v^2/c^2}},
\end{equation}
is precisely the relabeling that leaves $d\tau$ unchanged: it mixes time and space, tilting their axes toward
each other as the speed rises, while holding the interval fixed --- ours the tilt of the axes, the world's the interval they turn about ---
with $\gamma$ the factor by which clocks slow
and lengths contract as $v$ approaches $c$. The slowing is not a figure of speech but a measured fact. Muons
made high in the atmosphere by cosmic rays live, at rest, only about two microseconds --- far too short, even
at nearly the speed of light, to reach the ground before they decay. Yet they arrive in abundance, because at
their speed $\gamma$ is large and their clocks run slow: the few microseconds of their own proper time stretch,
in the ground's frame, into the tens of microseconds the descent takes. The muon reaches the ground because its
clock, not the Earth's, sets its lifetime.

The honest floor is a finite count obligation --- the interval is a
count of irreducible steps, invariant under any admissible relabeling. The reading is that there is no
preferred frame because there is no preferred labeling --- the labeling ours to assign, the interval-count the
world's to keep. What every observer agrees on is the interval, the count of steps along a history; the frames differ only in how they coordinate it, and the Lorentz
transformation is the relabeling that leaves the count alone.

\section{The Michelson--Morley Effect: the metric is isotropic}
\label{se:michelson-morley}

The Michelson--Morley experiment went looking for the ether wind and found nothing, and the null result is a
statement about the metric, not about a missing medium. The apparatus splits a beam of light along two
perpendicular arms, sends each to a mirror and back, and recombines them, watching the interference fringes for
any shift as the whole is slowly rotated --- the orientation ours to set, the fringes the world declines to
move. If light moved through an ether, the arm aligned with the Earth's
motion would take a different time than the crosswise arm, and turning the apparatus would sweep the fringes.
They never moved. In the language of the interval, the two arms are two equal-content paths, and that they
always agree --- in every orientation, at every season --- means the metric assigns the same interval to
equal-content paths regardless of direction: the axes ours to turn, the interval $\tau^2$ the world leaves
unchanged as they do.

The honest floor is a finite count obligation --- the count of steps along a path is isotropic, the same in
every direction. The reading is that the null result shows the metric isotropic --- the direction we set is
ours, the isotropy the world keeps in every one of them the world's. There is no ether to drag the light because
there is no preferred direction in the count; equal-content paths cost equal interval, and the
experiment's silence is the isotropy of the metric speaking.

\section{The Einstein Effect, I: the relativity of simultaneity}
\label{se:einstein-simultaneity}

The first of Einstein's two aspects is that simultaneity is not absolute. Look again at the boost,
\begin{equation}
  t' = \gamma\Bigl(t - \frac{vx}{c^2}\Bigr),
\end{equation}
and notice that the transformed time $t'$ depends not only on $t$ but on $x$: two events that happen at the same
time $t$ in one frame, but at different places, are assigned different times $t'$ in another. Simultaneity ---
the set of events sharing a clock reading --- is tilted by the boost; what is ``now'' for one observer is a
sloping slice of spacetime for another. The ``now'' we assign is ours to tilt with the frame; what can reach
what --- the causal order the interval fixes --- is the world's, the same in every frame. There is no
operational procedure that forces two observers in relative
motion to agree on the order of two spacelike-separated events: each one's refinement structure fixes its own
simultaneity, and no experiment can make them coincide. This is the relativity of distant order that Part II
met as the limit of synchronization --- no signal travels fast enough to impose a single ``now'' across space.

The honest floor is a finite count obligation --- simultaneity is a frame-relative partition of events, not an
absolute one, and no count forces agreement on the order of spacelike events. The reading is that simultaneity
is relative. The boost mixes time and space, so two observers slice ``now'' differently --- the slice ours to
assign, the causal order beneath it the world's to keep --- and there is no fact of the matter, beyond each
frame's own labeling, about which of two distant events came first.

\section{The Einstein Effect, II: proper time is maximal}
\label{se:einstein-proper-time}

The second aspect is about which path between two events a clock actually takes. Among all the admissible
worldlines connecting two events --- all the ways a clock could travel from the one to the other --- the path
physically realized is the one that maximizes the proper time,
\begin{equation}
  \tau = \int d\tau,
\end{equation}
the longest interval, the most distinguishable steps a clock can count between them --- the path ours to take,
the proper time the world scores it by the world's. This is the counterpart,
in spacetime, of the least-principles of the earlier parts: where a light ray took the path of least time and a
system the path of least action, a free clock takes the path of \emph{greatest} proper time. The famous
consequence is the twin paradox. Two twins part; one stays home on a straight worldline, the other accelerates
away and returns on a bent one; when they reunite, the traveler is younger. The straight worldline is the one
of maximal proper time --- counterintuitively, the longer path through space is the shorter path through time
--- so the stay-at-home twin, not the traveler, has aged more: the worldline is ours to run, and which twin the
world ages more --- the one on the straight path --- is the world's. The twin who travels the maximal-proper-time
path is the one who ages most.

The honest floor is a finite count obligation --- proper time is the maximal count of distinguishable steps
along an admissible chain. The reading is that the clock is not absolute but maximal. The proper time between
two events is the longest admissible chain of steps between them, and the realized worldline is the one that
counts the most --- the route ours to take, the tick-count the world keeps along it the world's --- the twin on
the straight, maximal path the one who ages most.

\section{The Refinement Effect: motion is differential refinement}
\label{se:refinement}

Relative motion is a difference in refinement depth. When one observer records more events than another over
the same stretch, no admissible refinement makes their event-orders agree; the mismatch is a difference in how
finely each resolves the causal order,
\begin{equation}
  \Delta N \text{ over } (a, b),
\end{equation}
and the observer recording more events is the one resolving the finer order --- the record ours to keep, the
causal order it resolves the world's. The honest floor is a finite count
obligation --- the difference is a count, $\Delta N$, of how many more events one record resolves than another.
The reading is that relative motion shows up as differential refinement. The observer who records more events
has the finer clock and resolves the finer causal order --- the depth each record reaches is ours, the one
history beneath both the world's --- so relative motion is not so much a thing moving through space as two
records resolving the same history to different depths.

\section{The Time Effect: the arrow is counted resolutions}
\label{se:time}

Time has an arrow because resolutions do not run backward. At each event the record carries a single unsatisfied
degree of freedom, and a tick of time is the discrete act of closing it --- turning an open refinement into a
closed record, eliminating one underdetermined degree of freedom and stabilizing what remains. Time is the
count of those closures,
\begin{equation}
  \text{time} = \#\{\text{closed degrees of freedom}\},
\end{equation}
each tick removing one --- the tick ours to mark, the one direction the count can go the world's. The honest
floor is a finite count obligation --- time is the count of resolved degrees
of freedom, one eliminated per tick. The reading is that the arrow of time is counted resolution. Each tick
closes an open degree of freedom and cannot reopen it, so the count only rises --- the clock ours to read, the
arrow beneath it the world's --- the asymmetry of time is the irreversibility of resolving the underdetermined,
the arrow pointing the way the count grows.

\section{LiDAR: the denser record ages more}
\label{se:lidar}

Time dilation can be measured by laser ranging. Two observers, idealized in flat space with gravity explicitly
set aside, exchange light pulses, and one of them accelerates; the return times become uneven in a way no single
inertial labeling can reconcile, the acceleration showing up as the asymmetry between the two paths. The clock
that accumulates the denser causal record --- more ticks, more resolved events --- has the larger proper time
and ages more --- the record ours to make dense or sparse by the path we take, the aging the world reads off its
count the world's. The effect has
been measured directly with clocks: in 1971 Hafele and Keating flew atomic clocks around the world aboard
commercial airliners and compared them against clocks left at rest, and the travelling clocks returned showing
just the offset relativity predicts, having ticked fewer times. The denser record is not a metaphor but a
reading on a caesium clock. The honest floor is
a smooth-shadow analogy --- the discrete refinement count is the model, the continuum Lorentz dilation its
smooth shadow, and the bridge from counted ticks to the dilation factor $\gamma$ is an analogy, named as one.
The reading is that the dilation is a difference in record density. Laser ranging reads it off the unevenness of
return times --- the ranging ours to run, the density it uncovers the world's --- and the twin with the denser
record is the one who has aged more --- proper time the count,
dilation its smooth shadow.

\section{The Sagnac Effect: rotation as a proved loop residue}
\label{se:sagnac}

The Sagnac effect is the third time the instrument closes a real loop and proves its residue. Send two light beams
around a closed ring, one clockwise and one counter-clockwise, and recombine them. If the ring is still, the two
travel equal paths and return in step. But set the ring rotating, and they no longer agree: the beam going with
the rotation must chase a receding finish line and takes longer, while the beam going against it meets an
approaching one and takes less, so the two return with a phase difference --- the ring ours to build and to
spin, the phase difference the rotation charges it with the world's. For a ring of area $A$ rotating at
angular velocity $\Omega$ the time difference is
\begin{equation}
  \Delta t = \frac{4\Omega A}{c^2},
\end{equation}
the Sagnac shift, proportional to the enclosed area and the rotation rate. Seen through the holonomy, this is
the residue of transporting timing around the loop: clocks carried around a rotating ring cannot be globally
synchronized, because the connection that adjusts local timing has a holonomy, and the forward and backward
legs return with a signed difference. Here, as in the echo maze and the Bragg peak --- the first and second
times the instrument closed a real loop and proved its residue --- the core is proved: the signed tally
difference between the co- and counter-rotating legs is a proved nonzero loop residue, a finite theorem about the
loop, not merely an analogy --- the loop ours to close, the residue the world will not let cancel.

The honest floor is a smooth-shadow analogy with a proved core --- the continuum $\Delta t = 4\Omega A/c^2$ is
the smooth shadow, but the nonzero signed residue beneath it is proved, the third such proof. The reading is
that rotation is a loop residue --- the ring ours to spin, the residue it cannot shed the world's. A rotating
frame cannot synchronize its clocks because the loop
carries a holonomy, and the finite model proves that residue nonzero --- the Sagnac shift the smooth shadow of a
proved failure to close, and the ring-laser gyroscope that holds an aircraft on course its working embodiment.

\bigskip

With frames, clocks, and the interval, Part V has its foundations. There is no preferred frame and no preferred
direction; the Minkowski interval is the one invariant every observer shares; simultaneity is relative and
proper time maximal; motion is a difference in refinement depth; the arrow is counted resolutions; and even
rotation is a proved loop residue --- the third of them. Part V turns next to gravity proper --- the
redshift, the delay, and the potentials that bend the record's own geometry.

\bigskip
\begin{center}
\rule{0.4\textwidth}{0.4pt}
\end{center}
\bigskip

\noindent The third witness takes the stand at last. Through nineteen chapters it was named and awaited --- the
reader that neither rides the wheel nor stands at the roadside, but looks down from far overhead and reads the
car by a clock of its own. Now it is here, and the first thing it says unsettles the court: there is no
privileged clock, and no privileged place from which to keep one. Every frame is a labeling, one among many, and
none of them the world's to prefer --- each reading the case has gathered was taken from somewhere, and no
somewhere is the true one.

This is the bending measure the last chapter promised, and it bends further than the court expected. Set two
clocks moving relative to one another and they will not keep the same time; ask two observers in motion which of
two distant events came first and they need not agree; what is one witness's single \emph{now} is another's
slanting slice of before-and-after. The overhead reader's clock runs at its own rate, neither the wheel's nor
the field's, so the three readings the case has gathered no longer even share a measure in which to be compared.
The frame is ours to choose; that none of them is privileged, and that the measures diverge, is the world's.
This is no figure of speech. Muons made high in the air, living at rest only a sliver of time, reach the ground
in numbers only because their own clocks, racing near the speed of light, run slow enough to stretch that sliver
across the whole descent; and atomic clocks flown around the world aboard airliners came home reading a little
behind the clocks left at rest, having ticked, in flight, a measurably smaller count. The bending is a reading
on a caesium clock, not a manner of speaking.

And yet the divergence is not chaos, for beneath the disagreeing clocks lies one thing every frame keeps the
same. Give any clock a path through events and it counts, along that path, a number all observers agree on
however they slice their time and their space: the interval, the proper time the clock itself tallies as it
goes. The count the book has carried since the first mark has, here, acquired a geometry --- and the geometry is
nothing but the measure of the count along its own path. Which frame reads it, and how the reading is cut into
time and space, is ours; the interval the clock counts is the world's, one and the same for all. Nor is any
direction privileged over another: light sent down two crossed arms and back returns in step however the whole
is turned, and that stubborn agreement is the measure declining to prefer a way to point. (Even a turning
leaves its mark: send a signal both ways around a spinning ring and the two do not come home in step --- a
proved failure to close, the third time the instrument carries a tally around a real loop and proves it comes
back charged, though what that residue is worth the book, as ever, keeps back.)

And the \emph{charge} takes on that geometry too. The charge was a count; now the count is measured along a path,
and no single standing clock reads it --- what a thing accumulates depends on the road it took to get there, and
two roads between the same two points need not tally the same. So the charge the case will collect is no longer
a number read off one dial; it is an interval counted along a path, the proper time of the thing that carried
it. The loop the charge counts is itself a path through events, and the interval is its true length. The twin who
stays home on a straight road ages more than the one who loops away and returns, because the straight path
counts the longer interval --- the most proper time between two meetings falls to the plainest route between
them. The trace
is that path-length held invariant --- the one measure that survives every witness's slicing, kept whole where
each frame's own clock would disagree. Whatever comes to be owed is measured along the road, not read from a
clock standing still.

And now, for the first time, all three witnesses are in the room together. The wheel's reading from within, the
field's from without, and the overhead reader's from above --- the three the case was built upon, present at
last, the meeting no longer a promise but a thing that can be assembled. It will not assemble itself. The third
witness reads by a clock that bends, and its measure does not fall into line with the other two of its own
accord; to bring the three onto one number, the case must first learn the correction that reconciles a bending
clock with standing ones. That reconciliation is the work the coming chapters take up --- the calibration that
lets three differently-clocked readings be made to meet. The three are on the bench; the number they must be
brought to is not yet in hand. The meeting has begun its work; it has not reached its end.

The overhead reader is only half-calibrated, though. What these chapters opened is the geometry of motion ---
clocks that slow with speed, frames with no privilege among them --- but the reader hangs high above the road,
and height, it will turn out, bends the measure as surely as speed does: a clock lifted up does not run at the
rate of one below. The next chapter turns to gravity proper --- the redshift, the delay, the pull that bends the
record's own geometry --- the other half of what the third witness needs before it can be believed. The case is
still open, its three witnesses finally assembled and not yet agreed; the court does not sit until the record is
complete.
